\documentclass[11pt, a4paper]{amsart}

\usepackage{amsmath, amssymb, amsthm}
\usepackage{enumerate}
\usepackage{hyperref}
\usepackage{booktabs}
\usepackage{array}
\usepackage{graphicx}
\usepackage[dvipsnames]{xcolor}

\theoremstyle{plain}
\newtheorem{theorem}{Theorem}[section]
\newtheorem{proposition}[theorem]{Proposition}
\newtheorem{lemma}[theorem]{Lemma}
\newtheorem{corollary}[theorem]{Corollary}
\newtheorem{conjecture}[theorem]{Conjecture}

\theoremstyle{definition}
\newtheorem{definition}[theorem]{Definition}
\newtheorem{remark}[theorem]{Remark}
\newtheorem{observation}[theorem]{Observation}

\newcommand{\ZZ}{\mathbb{Z}}
\newcommand{\RR}{\mathbb{R}}
\newcommand{\NN}{\mathbb{N}}
\renewcommand{\Pr}{\operatorname{Pr}}
\newcommand{\MI}{\operatorname{MI}}
\newcommand{\Ent}{\operatorname{H}}
\newcommand{\fS}{\mathfrak{S}}

\title[Information-Theoretic Properties of Prime Gap Sequences]{%
Information-Theoretic Properties of\\Prime Gap Sequences}

\author{Tobias Canavesi}
\address{Independent Researcher}
\email{tobiascanavesi@gmail.com}

\date{April 2026}

\begin{document}

\begin{abstract}
We study the prime gap sequence $g_n = p_{n+1} - p_n$ through the lens of information theory.
Our main result is a comparison of the empirical mutual information $\MI(g_n; g_{n+1}) \approx 0.33$ bits
against the Cram\'er random model, under which consecutive gaps are independent (so $\MI = 0$).
The full $0.33$ bits of lag-$1$ dependence is therefore attributable to the multiplicative structure of the integers---the sieve.
We verify this via permutation tests ($z > 1000$ at lag~$1$; lag~$2$ also significant; lags $\ge 3$ indistinguishable from noise) and
show that the non-Gaussian discrete character of prime gaps makes the mutual information approximately $300$ times larger than
the Gaussian prediction $\rho^2/(2\ln 2)$ from the linear autocorrelation $\rho_1 \approx -0.04$.
Supporting experiments include:
block entropy measurements with an explicit finite-sample ceiling analysis showing $H_k$ saturates at $\log_2 N$ for large $k$;
a Lempel--Ziv entropy rate estimator that avoids this curse of dimensionality, yielding $h_{\mathrm{LZ}}^{\mathrm{gaps}}/h_{\mathrm{LZ}}^{\mathrm{shuffled}} \approx 0.92$--$0.94$;
goodness-of-fit tests confirming monotonic convergence toward the Gallagher--Poisson prediction with variance-to-mean ratio climbing from $0.72$ to $0.77$;
and conditional entropy measurements revealing a ${\sim}0.06$-bit asymmetry between primes $\equiv 1$ and $\equiv 5 \pmod{6}$.
We also verify that the empirical single-gap entropy satisfies $H_1 \approx \log_2(\log x) + C$ with a slowly varying residual $C \approx 0.22$--$0.28$,
consistent with a discretized exponential model weighted by the Hardy--Littlewood singular series.
All code and data are publicly available.
\end{abstract}

\maketitle
\tableofcontents

% ====================================================================
\section{Introduction}
\label{sec:intro}
% ====================================================================

The prime gap sequence $g_n = p_{n+1} - p_n$ is one of the most studied objects in analytic number theory. Classical results characterize its average behavior (the prime number theorem gives $\langle g_n \rangle \sim \log p_n$), its extreme values (the Cram\'er conjecture~\cite{Cramer1936}, the Maynard--Tao theorem on bounded gaps~\cite{Maynard2015}), and its distributional properties (the Gallagher conjecture~\cite{Gallagher1976} predicting local Poisson count statistics conditional on the Hardy--Littlewood $k$-tuple conjecture). Lemke Oliver and Soundararajan~\cite{LemkeOliver2016} showed that consecutive primes exhibit unexpected biases in their joint distribution of residues modulo~$q$, traceable to the Hardy--Littlewood singular series.

Despite this extensive literature, the \emph{information-theoretic} properties of the gap sequence have received limited systematic attention.
Shannon entropy, mutual information, and algorithmic complexity provide a complementary lens that captures aspects of the sequence's structure---including nonlinear dependencies and compressibility---that linear statistical measures like autocorrelation cannot detect.
Cohen, Iyer, and Manack~\cite{Cohen2024} recently studied the exponential fit of normalized prime gaps, providing a distributional perspective that our work complements with information-theoretic tools.

The central contribution of this paper is a \emph{theoretical prediction}: under the Cram\'er random model, consecutive prime gaps are independent (Proposition~\ref{prop:cramer-mi}), so $\MI(g_n; g_{n+1}) = 0$.
The empirical value of $\approx 0.33$ bits, verified via permutation tests, is therefore entirely a signature of the multiplicative sieve.
We quantify this sieve contribution and compare it to predictions from the Hardy--Littlewood singular series.

\begin{remark}[Stationarity]
\label{rem:stationarity}
The prime gap sequence is not stationary: the gap distribution changes with scale as $\langle g_n \rangle \sim \log p_n$.
Throughout this paper, we treat each fixed-scale window of $N = 2 \times 10^5$ consecutive primes as approximately stationary.
All entropy-rate and mutual-information measurements are made within such windows; cross-scale comparisons reflect the evolution of the local gap distribution, not a single stationary process.
\end{remark}

% ====================================================================
\section{Definitions and Setup}
\label{sec:setup}
% ====================================================================

Let $p_1 = 2 < p_2 = 3 < p_3 = 5 < \cdots$ be the primes. Fix a starting scale $x$ and let $p_1(x) < p_2(x) < \cdots < p_N(x)$ be $N$ consecutive primes starting near $x$. Define the gap sequence $g_i = p_{i+1}(x) - p_i(x)$ for $i = 1, \ldots, N-1$.

\begin{definition}[Block entropy]
For $k \ge 1$, define the $k$-block entropy
\[
H_k = - \sum_{\mathbf{w} \in \mathcal{A}^k} \hat{p}(\mathbf{w}) \log_2 \hat{p}(\mathbf{w}),
\]
where $\hat{p}(\mathbf{w})$ is the empirical frequency of the $k$-tuple $\mathbf{w} = (w_1, \ldots, w_k)$ among the overlapping windows $(g_i, g_{i+1}, \ldots, g_{i+k-1})$, and $\mathcal{A} = \{2, 4, 6, 8, \ldots\}$ is the alphabet of even positive integers.
\end{definition}

\begin{definition}[Entropy rate]
The $k$-th entropy rate estimate is $h_k = H_k - H_{k-1}$ for $k \ge 2$, with $h_1 = H_1$. For a stationary ergodic process, $h_k \to h$ as $k \to \infty$, where $h$ is the entropy rate of the process.
\end{definition}

For an i.i.d.\ process, $H_k = k \cdot H_1$ and $h_k = H_1$ for all $k$. The \emph{redundancy} $R_k = k \cdot H_1 - H_k \ge 0$ measures the total dependence captured up to block size $k$. Note that $R_k$ equals the \emph{total correlation} (multi-information) of the $k$-tuple $(g_n, \ldots, g_{n+k-1})$.

\begin{definition}[Mutual information]
For random variables $X, Y$ with discrete joint distribution,
\[
\MI(X; Y) = H(X) + H(Y) - H(X, Y).
\]
We apply the Miller--Madow bias correction~\cite{Miller1955}: subtract $(|\mathcal{A}_{XY}| - |\mathcal{A}_X| - |\mathcal{A}_Y| + 1) / (2 n \ln 2)$ from the raw plug-in estimate, where $|\mathcal{A}_{XY}|$ is the number of observed joint categories.
\end{definition}

\begin{definition}[Lempel--Ziv complexity and entropy rate]
The LZ76 complexity $C(s)$ of a finite sequence $s = s_1 s_2 \cdots s_n$ is the number of distinct phrases in the sequential parsing of $s$ into shortest substrings not previously encountered~\cite{LZ76}.
For a stationary ergodic process over an alphabet of size $|\mathcal{A}|$,
\[
\frac{C(s) \cdot \log_{|\mathcal{A}|} n}{n} \to h \quad \text{almost surely as } n \to \infty,
\]
where $h$ is the entropy rate~\cite{LZ76}. We use this as a \emph{consistent} entropy rate estimator that avoids the curse of dimensionality inherent in plug-in block entropy for large $k$.
\end{definition}


% ====================================================================
\section{The Cram\'er Model and Mutual Information}
\label{sec:cramer}
% ====================================================================

We begin with the paper's main theoretical result: a comparison of the empirical lag-$1$ mutual information to the prediction of the Cram\'er random model.

\begin{proposition}[Independence under the Cram\'er model]
\label{prop:cramer-mi}
In the Cram\'er random model near $x$ --- where each integer $m > p$ is independently ``prime'' with probability $1/\log x$ --- consecutive gaps $g_n$ and $g_{n+1}$ are independent. Consequently, $\MI(g_n; g_{n+1}) = 0$ under this model.
\end{proposition}

\begin{proof}
In the constant-density Cram\'er model, each integer is independently prime with probability $1/\log x$, so gaps are i.i.d.\ $\mathrm{Geometric}(1/\log x)$.
Independence of consecutive gaps follows from the memoryless property: the waiting time from $p_{n+1}$ to $p_{n+2}$ is independent of the waiting time from $p_n$ to $p_{n+1}$, since the geometric distribution is the unique discrete memoryless distribution.
\end{proof}

\begin{remark}
In the Cram\'er model with varying density ($1/\log m$ depending on $m$), there is a residual coupling of order $O(1/(\log x)^2)$ because a large $g_n$ shifts $p_{n+1}$ to a region where $1/\log m$ is slightly different. This coupling is negligible at all scales we consider.
\end{remark}

\begin{observation}[The sieve contribution to MI]
\label{obs:sieve-mi}
Since $\MI_{\text{Cram\'er}} = 0$ and $\MI_{\text{empirical}} \approx 0.33$ bits at lag~$1$, the entire observed mutual information is attributable to the multiplicative structure of the integers---the sieve.

Under the Hardy--Littlewood conjecture, the probability of a gap of size $h$ near $x$ is
\begin{equation}\label{eq:hl-marginal}
\Pr(g_n = h) \approx \frac{\fS(\{0, h\})}{\log x} \, e^{-h/\log x},
\end{equation}
where $\fS(\{0, h\}) = 2C_2 \prod_{p \mid h,\, p > 2} (p-1)/(p-2)$ depends on the prime factorization of $h$.
For \emph{consecutive} gap pairs, the joint probability involves three primes $(p_n, p_n + h_1, p_n + h_1 + h_2)$ and is governed by the triple singular series:
\begin{equation}\label{eq:hl-joint}
\Pr(g_n = h_1,\, g_{n+1} = h_2) \approx \frac{\fS(\{0, h_1, h_1 + h_2\})}{(\log x)^2} \, e^{-(h_1 + h_2)/\log x}.
\end{equation}
Mutual information arises because the joint~\eqref{eq:hl-joint} does not factor into the product of marginals~\eqref{eq:hl-marginal}: in general, $\fS(\{0, h_1, h_1 + h_2\}) \ne \fS(\{0, h_1\}) \cdot \fS(\{0, h_2\}) / (2 C_2)$ due to shared small-prime divisibility constraints among $h_1$, $h_2$, and $h_1 + h_2$.
\end{observation}

\subsection{Quantitative prediction under Hardy--Littlewood}

We compute $\MI_{\text{H-L}}$ numerically by summing~\eqref{eq:hl-marginal} and~\eqref{eq:hl-joint} over all gap pairs $(h_1, h_2)$ up to $\sim 8 \log x$ (the result is insensitive to this cutoff due to the rapid decay of $e^{-(h_1+h_2)/\log x}$), computing the marginal and joint entropies, and taking their difference.
For the triple singular series appearing in~\eqref{eq:hl-joint}, we use the product formula
\[
\fS(\{0, h_1, h_1 + h_2\}) = \prod_{p} \frac{(1 - \nu_p / p)}{(1 - 1/p)^3},
\]
where $\nu_p = |\{0, h_1, h_1 + h_2\} \bmod p|$ counts distinct residues.
The product is taken over all primes $p$; since $\nu_p = 3$ for all $p > h_1 + h_2$ (the three elements are automatically distinct modulo large primes), only finitely many factors differ from the baseline, and the product converges rapidly.

\begin{table}[ht]
\centering
\caption{Mutual information: Cram\'er model, Hardy--Littlewood prediction, and empirical.}
\label{tab:cramer-mi}
\begin{tabular}{ccccccc}
\toprule
Scale & $\log x$ & $\MI_{\text{Cram\'er}}$ & $\MI_{\text{H-L}}$ & $\MI_{\text{emp.}}$ & Residual & H-L / emp. \\
\midrule
$10^4$  & 14.1 & 0 & 0.239 & 0.339 & 0.100 & 70.5\% \\
$10^5$  & 14.2 & 0 & 0.239 & 0.338 & 0.099 & 70.7\% \\
$10^6$  & 14.7 & 0 & 0.238 & 0.335 & 0.098 & 70.9\% \\
$10^7$  & 16.3 & 0 & 0.234 & 0.329 & 0.095 & 71.1\% \\
$10^8$  & 18.4 & 0 & 0.229 & 0.322 & 0.093 & 71.2\% \\
\bottomrule
\end{tabular}
\end{table}

\begin{observation}[Quantitative decomposition]
The Hardy--Littlewood model, computed using the triple singular series $\fS(\{0, h_1, h_1 + h_2\})$ for all prime factors up to $\max(h_1, h_1+h_2)$, captures $\approx 71\%$ of the empirical MI, with remarkable stability across four orders of magnitude in scale.
The residual ($\approx 0.09$--$0.10$ bits) slowly decreases with scale and likely reflects finite-range effects: the asymptotic Hardy--Littlewood prediction~\eqref{eq:hl-joint} omits error terms of order $O(1/(\log x)^A)$ and the exponential weighting $e^{-(h_1+h_2)/\log x}$ is itself an approximation.
That the H-L prediction accounts for a constant fraction of the MI, independent of scale, suggests the singular series captures the dominant mechanism and the residual will vanish as $x \to \infty$.
\end{observation}


% ====================================================================
\section{Block Entropy and the Finite-Sample Ceiling}
\label{sec:block-entropy}
% ====================================================================

Table~\ref{tab:block-entropy} summarizes the block entropy measurements at the largest scale ($10^8$).

\begin{table}[ht]
\centering
\caption{Block entropy at scale $10^8$ ($N = 200{,}000$ primes).}
\label{tab:block-entropy}
\begin{tabular}{ccccc}
\toprule
$k$ & $H_k$ (bits) & $h_k$ (bits) & $H_k^{\mathrm{iid}}$ (bits) & $n_{\text{distinct}}$ \\
\midrule
1  & 4.481 & 4.481 & 4.481  & 80 \\
2  & 8.634 & 4.153 & 8.963  & 1{,}879 \\
3  & 12.499 & 3.865 & 13.444 & 18{,}671 \\
4  & 15.607 & 3.109 & 17.926 & 81{,}651 \\
5  & 17.186 & 1.579 & 22.407 & 163{,}726 \\
6  & 17.557 & 0.370 & 26.889 & 194{,}835 \\
7  & 17.604 & 0.047 & 31.370 & 199{,}431 \\
8  & 17.609 & 0.005 & 35.852 & 199{,}929 \\
9  & 17.610 & 0.001 & 40.333 & 199{,}986 \\
10 & 17.610 & 0.000 & 44.815 & 199{,}990 \\
\bottomrule
\end{tabular}
\end{table}

\begin{remark}[Finite-sample ceiling]
\label{rem:ceiling}
The saturation of $H_k$ near $\log_2 N \approx 17.6$ bits is primarily a \emph{finite-sample artifact}, not an intrinsic property of the gap process. At $k = 10$, nearly every overlapping window is unique ($n_{\text{distinct}} \approx N$), so the plug-in estimator returns $H_{10} \approx \log_2 N$ regardless of the true entropy rate.
The entropy rate $h_k$ appearing to converge to zero for large $k$ reflects this ceiling, not a genuine vanishing of conditional uncertainty.
The true entropy rate of the gap sequence (within a fixed-scale window) is almost certainly positive.
\end{remark}

The block entropy estimates are \emph{reliable} for small $k$ where $n_{\text{distinct}} \ll N$. In particular:
\begin{itemize}
\item At $k = 1$: $n_{\text{distinct}} = 80 \ll 199{,}999$. The estimate $H_1 = 4.481$ bits is well-determined.
\item At $k = 2$: $n_{\text{distinct}} = 1{,}879 \ll 199{,}998$. The estimate $H_2 = 8.634$ bits and the drop $h_2 = 4.153 < H_1$ are meaningful, revealing genuine lag-$1$ dependence.
\item At $k \ge 5$: $n_{\text{distinct}}$ approaches $N$, and the estimates begin to be dominated by the sample-size ceiling.
\end{itemize}

The redundancy $R_2 = 2 H_1 - H_2 = 0.33$ bits at $k = 2$ is a robust measurement that matches the lag-$1$ mutual information (as it must: $R_2 = \MI(g_n; g_{n+1})$).


\subsection{The LZ entropy rate estimator}

To obtain a consistent entropy rate estimate free of the curse of dimensionality, we use the LZ76 complexity.
The ratio $h_{\mathrm{LZ}}^{\mathrm{gaps}} / h_{\mathrm{LZ}}^{\mathrm{shuffled}}$ measures the fraction of marginal entropy that survives as true process entropy:

\begin{table}[ht]
\centering
\caption{LZ-based entropy rate estimates. $h_{\mathrm{LZ}} = C \cdot \log_{|\mathcal{A}|}(n) / n$.}
\label{tab:lz-rate}
\begin{tabular}{cccc}
\toprule
Scale & $h_{\mathrm{LZ}}^{\mathrm{gaps}}$ & $h_{\mathrm{LZ}}^{\mathrm{shuffled}}$ & Ratio \\
\midrule
$10^4$ & 0.6538 & 0.7079 & 0.9235 \\
$10^5$ & 0.6580 & 0.7105 & 0.9260 \\
$10^6$ & 0.6644 & 0.7152 & 0.9290 \\
$10^7$ & 0.6685 & 0.7143 & 0.9359 \\
$10^8$ & 0.6825 & 0.7237 & 0.9431 \\
\bottomrule
\end{tabular}
\end{table}

The ratio $\approx 0.92$--$0.94$ is the key quantity: the gap sequence carries about $6$--$8\%$ less entropy per symbol than an i.i.d.\ process with the same marginal distribution. This monotonically increases toward~$1$ with scale, consistent with the Gallagher--Poisson prediction.

\subsection{Scaling of $H_1$}

Under a discretized exponential model weighted by the singular series, the single-gap entropy should scale as
\begin{equation}\label{eq:H1-scaling}
H_1 \approx \log_2(\log x) + C
\end{equation}
for a constant $C$ that depends on the singular series weights. Empirically:

\begin{table}[ht]
\centering
\caption{$H_1$ vs.\ $\log_2(\log x)$.}
\label{tab:h1-scaling}
\begin{tabular}{cccc}
\toprule
Scale & $H_1$ (bits) & $\log_2(\log x)$ & Residual \\
\midrule
$10^4$ & 4.035 & 3.817 & 0.218 \\
$10^5$ & 4.053 & 3.825 & 0.228 \\
$10^6$ & 4.128 & 3.878 & 0.250 \\
$10^7$ & 4.289 & 4.024 & 0.265 \\
$10^8$ & 4.481 & 4.205 & 0.277 \\
\bottomrule
\end{tabular}
\end{table}

The leading term $\log_2(\log x)$ captures the bulk of the growth. The residual \emph{increases} slowly with scale (from $0.22$ to $0.28$), reflecting the broadening gap alphabet: at larger scales, rarer large gaps appear, adding marginal entropy beyond the continuous exponential prediction. Deriving a precise asymptotic for this residual from the singular series remains an open problem (see \S\ref{sec:discussion}).


% ====================================================================
\section{Mutual Information: Empirical Analysis}
\label{sec:mi}
% ====================================================================

\subsection{Lag-1 MI and the Gaussian comparison}

The mutual information $\MI(g_n; g_{n+\ell})$ captures all statistical dependence between gaps separated by $\ell$ positions.
At lag~$1$, the empirical MI is $\approx 0.32$--$0.34$ bits across all scales. For comparison, the Gaussian approximation gives $\MI \approx \rho^2/(2 \ln 2)$ for bivariate Gaussian variables with correlation $\rho$. With $\rho_1 \approx -0.04$, this predicts $\MI \approx 10^{-3}$ bits.

\begin{observation}
\label{obs:nonlinear}
The discrepancy of a factor $\approx 300$ between $\MI_{\text{empirical}}$ and the Gaussian prediction reflects the non-Gaussian, discrete character of the gap distribution---not a surprising failure of a linear model, but rather a quantification of how much information is carried by the full discrete joint distribution $(g_n, g_{n+1})$ beyond its second moments. The gap distribution is supported on $\{2, 4, 6, \ldots\}$ with a long right tail, making the Gaussian approximation inappropriate.
\end{observation}

\subsection{Lag $\ell \ge 2$ and permutation tests}

At lags $\ell \ge 2$, the Miller--Madow-corrected MI values are 0.002--0.005 bits, which is on the same order as the bias correction itself ($\approx 0.003$--$0.005$ bits for an alphabet of $\sim 30$--$80$ values). To determine whether these values represent genuine dependence or estimator noise, we perform permutation tests: for each lag $\ell$, we shuffle the sequence $(g_1, \ldots, g_{n-\ell})$ while keeping $(g_{1+\ell}, \ldots, g_n)$ fixed, compute MI on $200$ such random shuffles, and compare the empirical MI to this null distribution. This preserves both marginal distributions while destroying sequential structure.

\begin{table}[ht]
\centering
\caption{Permutation test results for MI significance at scale $10^8$ (200 shuffles).}
\label{tab:perm}
\begin{tabular}{ccccc}
\toprule
Lag & MI observed (bits) & Null mean (bits) & $z$-score & $p$-value \\
\midrule
1 & 0.3286 & 0.0122 & 1148 & $< 0.005$ \\
2 & 0.0149 & 0.0122 & 11.2 & $< 0.005$ \\
3 & 0.0116 & 0.0122 & $-1.9$ & 0.97 \\
5 & 0.0121 & 0.0121 & $-0.03$ & 0.53 \\
10 & 0.0120 & 0.0121 & $-0.6$ & 0.74 \\
50 & 0.0119 & 0.0121 & $-1.1$ & 0.90 \\
100 & 0.0123 & 0.0122 & $0.3$ & 0.38 \\
\bottomrule
\end{tabular}
\end{table}

The permutation tests yield a critical finding: lag-$1$ MI is overwhelmingly significant ($z > 1000$), and lag-$2$ MI is also significant ($z \approx 11$), but \textbf{at all lags $\ell \ge 3$, the MI is indistinguishable from the permutation null} ($|z| < 2$, $p > 0.01$). The Miller--Madow-corrected MI values reported for $\ell \ge 3$ in Experiment~2 are within the estimator's noise floor and should not be interpreted as evidence of genuine dependence. The gap process appears to have \emph{effective memory of length~$2$}: $g_n$ and $g_{n+1}$ share the prime $p_{n+1}$, and $g_n$ and $g_{n+2}$ share a weak residual correlation (possibly mediated by modular structure), but beyond lag~$2$ the process is effectively independent.

\begin{remark}
The permutation test null distribution measures the MI expected from \emph{marginal} structure alone (same gap frequencies, destroyed sequential order). Any MI exceeding the null reflects sequential structure in the gap process.
\end{remark}


% ====================================================================
\section{Lempel--Ziv Complexity}
\label{sec:lz}
% ====================================================================

Table~\ref{tab:lz} reports the LZ76 complexity of the gap sequence compared to shuffled and Poisson surrogates.

\begin{table}[ht]
\centering
\caption{LZ76 complexity at each scale ($n = 199{,}999$ gaps, 20 surrogates).}
\label{tab:lz}
\begin{tabular}{cccccc}
\toprule
Scale & $C_{\text{gaps}}$ & $C_{\text{shuffled}}$ & $C_{\text{Poisson}}$ & $C/C_{\text{shuf}}$ & $C/C_{\text{Poi}}$ \\
\midrule
$10^4$ & 44{,}212 & 47{,}871 & 48{,}472 & 0.924 & 0.912 \\
$10^5$ & 44{,}494 & 48{,}035 & 48{,}655 & 0.926 & 0.914 \\
$10^6$ & 45{,}278 & 48{,}736 & 49{,}337 & 0.929 & 0.918 \\
$10^7$ & 46{,}999 & 50{,}222 & 50{,}829 & 0.936 & 0.925 \\
$10^8$ & 49{,}007 & 51{,}960 & 52{,}546 & 0.943 & 0.933 \\
\bottomrule
\end{tabular}
\end{table}

The gap between $C_{\text{gaps}}$ and $C_{\text{shuffled}}$ reflects sequential structure destroyed by shuffling. Since $C_{\text{gaps}} / C_{\text{shuffled}}$ equals the ratio of LZ entropy rates (see Table~\ref{tab:lz-rate}), this provides a dimensionality-free confirmation that the gap sequence carries about $6$--$8\%$ less entropy than its marginal distribution alone would suggest.
The ratio approaches~$1$ from below, supporting the Gallagher prediction qualitatively.


% ====================================================================
\section{Goodness-of-Fit to the Gallagher--Poisson Model}
\label{sec:gof}
% ====================================================================

Gallagher~\cite{Gallagher1976} showed, conditionally on the Hardy--Littlewood $k$-tuple conjecture, that the number of primes in a random interval of length $\lambda \log x$ near $x$, averaged over starting points $n \le N$, follows a Poisson distribution with mean $\lambda$.
For a Poisson process, the inter-arrival times are i.i.d.\ exponential, so this implies (under standard Poisson process theory) that the normalized gaps $g_n / \langle g_n \rangle$ should converge in distribution to Exp$(1)$.
We note that Gallagher's theorem is stated in terms of counts averaged over starting points, not directly as a statement about the gap sequence; the exponential convergence is a standard inference from Poisson count statistics~\cite{Cohen2024}.

\subsection{Exponential fit}

\begin{table}[ht]
\centering
\caption{Goodness-of-fit statistics for $g_n / \langle g_n \rangle$ against Exp$(1)$.}
\label{tab:exp-fit}
\begin{tabular}{cccc}
\toprule
Scale & KS statistic & AD statistic & $\chi^2$ statistic \\
\midrule
$10^4$ & 0.1597 & 5{,}463 & 147{,}349 \\
$10^5$ & 0.1592 & 5{,}406 & 145{,}143 \\
$10^6$ & 0.1541 & 5{,}007 & 145{,}205 \\
$10^7$ & 0.1466 & 4{,}445 & 107{,}294 \\
$10^8$ & 0.1343 & 3{,}659 & 82{,}103 \\
\bottomrule
\end{tabular}
\end{table}

All tests reject at every scale, which is expected given $n \approx 2 \times 10^5$. However, all three test statistics decrease monotonically with scale, confirming convergence.

\subsection{Poisson counting and underdispersion}

\begin{table}[ht]
\centering
\caption{Variance-to-mean ratio for prime counts in Poisson windows ($\lambda = 1$). Poisson predicts $\sigma^2/\mu = 1$.}
\label{tab:poisson}
\begin{tabular}{cc}
\toprule
Scale & $\sigma^2 / \mu$ \\
\midrule
$10^4$ & 0.722 \\
$10^5$ & 0.723 \\
$10^6$ & 0.733 \\
$10^7$ & 0.746 \\
$10^8$ & 0.768 \\
\bottomrule
\end{tabular}
\end{table}

\begin{observation}[Underdispersion]
Prime counts are systematically \emph{underdispersed} relative to Poisson ($\sigma^2/\mu < 1$) at all scales.
This reflects the sieve's ``repulsion'' effect: the presence of a prime at $p$ eliminates multiples of $p$ from nearby candidacy, reducing the variance of prime counts below the Poisson level.
The ratio increases monotonically toward~$1$, consistent with convergence to the Gallagher prediction.
\end{observation}


% ====================================================================
\section{Conditional Entropy by Residue Class}
\label{sec:conditional}
% ====================================================================

All primes $p > 3$ satisfy $p \equiv 1$ or $5 \pmod{6}$. We partition primes by residue class and measure the gap entropy within each class.

\begin{table}[ht]
\centering
\caption{Single-gap entropy $H_1$ (bits) by prime residue class mod~$6$.}
\label{tab:cond-entropy}
\begin{tabular}{cccc}
\toprule
Scale & $H_1$ (all) & $H_1$ ($p \equiv 1$) & $H_1$ ($p \equiv 5$) \\
\midrule
$10^4$ & 4.035 & 6.882 & 6.924 \\
$10^5$ & 4.053 & 6.921 & 6.964 \\
$10^6$ & 4.128 & 7.073 & 7.122 \\
$10^7$ & 4.289 & 7.393 & 7.468 \\
$10^8$ & 4.481 & 7.799 & 7.860 \\
\bottomrule
\end{tabular}
\end{table}

\begin{observation}
Primes $\equiv 5 \pmod{6}$ have consistently higher gap entropy ($\approx 0.05$--$0.07$ bits) than primes $\equiv 1 \pmod{6}$. We conjecture that this traces to the Hardy--Littlewood singular series: for the pair $\{0, h\}$, $\fS(\{0, h\}) = 2C_2 \prod_{p \mid h,\, p > 2} (p-1)/(p-2)$, so gaps divisible by~$3$ (like $g = 6$) are enhanced by a factor $(3-1)/(3-2) = 2$ relative to gaps not divisible by~$3$ (like $g = 4$).
Since the allowed gap values and their singular-series weights differ between residue classes (gaps between primes $\equiv 1 \pmod{6}$ must be $\equiv 0 \pmod{6}$, and similarly for class $5$), the resulting gap distributions have slightly different entropies. A precise derivation of the predicted entropy difference from the singular series is an interesting open problem.
\end{observation}

\begin{remark}
The class-restricted entropies ($\approx 7$--$8$ bits) are much larger than the full-sequence entropy ($\approx 4$--$4.5$ bits) because class-restricted gaps span roughly twice the range of all-prime gaps, yielding a broader distribution with more entropy.
\end{remark}


% ====================================================================
\section{Discussion}
\label{sec:discussion}
% ====================================================================

\subsection{The sieve as the source of dependence}

The central result of this paper is that consecutive prime gaps carry $\approx 0.33$ bits of mutual information, and that this dependence is entirely absent from the Cram\'er model (Proposition~\ref{prop:cramer-mi}).
The sieve creates this dependence through the singular series: the joint probability of a specific gap pair $(g_n, g_{n+1})$ depends on the factorization structure of the triple $(0, g_n, g_n + g_{n+1})$, which generally does not decouple.
The non-Gaussian character of the discrete gap distribution amplifies the MI far beyond what the modest linear autocorrelation would suggest.

\subsection{Convergence to the Gallagher prediction}

Five independent measurements confirm that the gap sequence becomes ``more random'' at larger scales:
\begin{itemize}
\item LZ entropy rate ratio increases from $0.924$ to $0.943$ (Table~\ref{tab:lz-rate}).
\item MI at lag~$1$ decreases from $0.34$ to $0.32$ bits.
\item KS and AD test statistics decrease monotonically (Table~\ref{tab:exp-fit}).
\item Variance-to-mean ratio increases from $0.72$ to $0.77$ (Table~\ref{tab:poisson}).
\item The residual in $H_1 - \log_2(\log x)$ increases slowly from $0.22$ to $0.28$ (Table~\ref{tab:h1-scaling}), reflecting the broadening gap alphabet.
\end{itemize}

These provide independent quantifications of the \emph{rate} of convergence to the Poisson/exponential regime. The convergence is real but slow.

\subsection{Open directions}

\begin{enumerate}
\item \textbf{Asymptotics of $\MI_{\text{H-L}}$}: our numerical computation (Table~\ref{tab:cramer-mi}) shows $\MI_{\text{H-L}} \approx 0.23$ bits, capturing $71\%$ of the empirical MI with remarkable scale-stability. Prove that $\MI_{\text{H-L}}(x) \to 0$ as $x \to \infty$ (as the Gallagher prediction requires) and determine the rate. Prove or disprove that the ratio $\MI_{\text{H-L}} / \MI_{\text{empirical}}$ has a limit as $x \to \infty$.

\item \textbf{Prove $H_1 = \log_2(\log x) + C + o(1)$}: under the Hardy--Littlewood conjecture, derive the constant $C$ from the singular series weights and the discretization to even integers.

\item \textbf{Multi-information decay}: the redundancy $R_k = k H_1 - H_k$ (equivalently, the total correlation) as a function of $k$ should, for a Markov process of order $m$, grow linearly for $k \le m$ and then saturate. The shape of $R_k$ at small $k$ (where estimates are reliable) could determine the effective memory length of the gap process.
\end{enumerate}


% ====================================================================
\section{Computational Details}
\label{sec:computation}
% ====================================================================

All computations were performed in Python~3 using SymPy~\cite{SymPy} for prime generation, NumPy for array operations, and SciPy for statistical tests. Permutation tests used $200$ shuffles per lag. The total computation time was approximately 10~minutes on a standard laptop. All code and data are available at:

\begin{center}
\url{https://github.com/tobiascanavesi/information-theory-gaps}
\end{center}

\subsection*{Use of generative AI}
The generative AI tool Claude Opus 4.6 (Anthropic, 2025) was used for assistance with Python code development and \LaTeX{} formatting. All mathematical content and interpretation represent the original intellectual contribution of the author.


% ====================================================================
% References
% ====================================================================

\begin{thebibliography}{99}

\bibitem{Cramer1936}
H.~Cram\'er, \emph{On the order of magnitude of the difference between consecutive prime numbers}, Acta Arithmetica \textbf{2} (1936), 23--46.

\bibitem{Gallagher1976}
P.~X.~Gallagher, \emph{On the distribution of primes in short intervals}, Mathematika \textbf{23} (1976), 4--9.

\bibitem{Maynard2015}
J.~Maynard, \emph{Small gaps between primes}, Annals of Mathematics \textbf{181} (2015), no.~1, 383--413.

\bibitem{LemkeOliver2016}
R.~J.~Lemke Oliver and K.~Soundararajan, \emph{Unexpected biases in the distribution of consecutive primes}, Proc.\ Natl.\ Acad.\ Sci.\ USA \textbf{113} (2016), no.~31, E4446--E4454.

\bibitem{Cohen2024}
A.~Cohen, S.~Iyer, and C.~Manack, \emph{Gaps between consecutive primes and the exponential distribution}, Experimental Mathematics (2024).

\bibitem{LZ76}
A.~Lempel and J.~Ziv, \emph{On the complexity of finite sequences}, IEEE Trans.\ Inform.\ Theory \textbf{22} (1976), no.~1, 75--81.

\bibitem{Miller1955}
G.~A.~Miller, \emph{Note on the bias of information estimates}, in: H.~Quastler (ed.), Information Theory in Psychology: Problems and Methods, Free Press, Glencoe, IL, 1955, pp.~95--100.

\bibitem{SymPy}
A.~Meurer et al., \emph{SymPy: symbolic mathematics in Python}, PeerJ Computer Science \textbf{3} (2017), e103.

\end{thebibliography}

\end{document}
